\documentclass[11pt]{article}
\usepackage[letterpaper,margin=0.9in]{geometry}
\usepackage{amsmath,amssymb,amsthm}
\usepackage{pgfplots}
\pgfplotsset{compat=1.18}
\usetikzlibrary{positioning}
\usepackage{tcolorbox}
\usepackage{listings}
\usepackage{booktabs}
\usepackage[colorlinks,allcolors=blue!50!black]{hyperref}
\usepackage{needspace}
% LaTeX's defaults; Debian's TeX Live 2023 leaves these at 0 ("u-nique")
\lefthyphenmin=2 \righthyphenmin=3
% stretch a tight paragraph instead of overflowing (two-digit citations)
\emergencystretch=1em
% start a section on a new page if less than a quarter page remains
\let\oldsection\section
\renewcommand{\section}{\Needspace{0.25\textheight}\oldsection}

\newtheorem{theorem}{Theorem}[section]
\newtheorem{proposition}[theorem]{Proposition}
\newtheorem{corollary}[theorem]{Corollary}
\theoremstyle{definition}
\newtheorem{definition}[theorem]{Definition}
\newtheorem{example}[theorem]{Example}
\theoremstyle{remark}
\newtheorem{remark}[theorem]{Remark}

\lstset{
  basicstyle=\ttfamily\small,
  columns=fullflexible,
  keepspaces=true,
  keywordstyle=\color{blue!50!black},
  commentstyle=\color{black!60}\itshape,
  captionpos=t,
  aboveskip=0.8em,
  belowskip=0.4em,
}

\title{
  \textbf{Functional Versus Imperative Is Solved}\\[0.2em]
  \large Power, Cost, and Shared State in the Paradigm Debate
}
\author{}
\date{29 October 2026}

\begin{document}
\maketitle

\section{Two Ways to Sum a List}\label{sec:defs}

Here is the same function written twice.

\noindent
\begin{minipage}[t]{0.5\textwidth}
  \begin{lstlisting}[language=C, caption={Imperative, in C.}, label=lst:c]
int total(const int *xs, int n) {
  int acc = 0;
  for (int i = 0; i < n; i++)
    acc += xs[i];
  return acc;
}
  \end{lstlisting}
\end{minipage}%
\begin{minipage}[t]{0.5\textwidth}
  \begin{lstlisting}[language=Haskell, caption={Functional, in Haskell.}, label=lst:hs]
total :: [Int] -> Int
total []       = 0
total (x : xs) = x + total xs
  \end{lstlisting}
\end{minipage}

\medskip
\noindent The first is a list of instructions. The second is a pair of
equations. Programmers have argued about which is better for most of the
history of programming, and the argument is older than most of the
languages it is now conducted in. We settle it.

\begin{definition}[Imperative]
  A program is \emph{imperative} if its meaning is a sequence of
  commands, each of which changes a \emph{store}, a map from locations to
  values. Running it executes the commands in order, and its result is
  read from the final store.
\end{definition}

\begin{definition}[Functional]
  A program is \emph{functional} if its meaning is an expression, and the
  meaning of an expression is its value. It is \emph{pure} if the value
  of every expression depends only on the values of its subexpressions,
  so that any subexpression can be replaced by its value without changing
  the meaning of the whole. Strachey, borrowing from Quine, called this
  property \emph{referential transparency}~\cite{strachey}.
\end{definition}

\begin{remark}
  The term has since been used in several senses that are not
  equivalent~\cite{sondergaard}. We use Strachey's. Both definitions
  classify programs, not languages. C permits a function that assigns
  nothing, and Haskell permits one that does nothing but assign. Most
  languages in use permit both, so the question is not which language to
  use but which kind of program to write.
\end{remark}

\section{The Ceasefire of 1937}\label{sec:power}

The first question, whether one style can compute something the other
cannot, was settled before either had a machine to run on. Church
proposed that the effectively calculable functions are those definable
in the $\lambda$-calculus~\cite{church}. Turing proposed those computable
by a machine that reads and writes symbols on a tape, following a table
of instructions~\cite{turing36}. The first is a functional language and
the second an imperative one.

\begin{theorem}[The Church--Turing Ceasefire]\label{thm:ceasefire}
  A function on the natural numbers is $\lambda$-definable if and only if
  it is computable by a Turing machine.
\end{theorem}

\begin{proof}
  Turing showed that every $\lambda$-definable function is computable and
  every computable function is general recursive~\cite{turing37}, and
  Kleene that every general recursive function is
  $\lambda$-definable~\cite{kleene}.
\end{proof}

\begin{corollary}[No Exclusive Programs]\label{cor:exclusive}
  Any function computed by a program in one Turing-complete language,
  functional or imperative, is computed by some program in any other.
\end{corollary}

\begin{proof}
  Each such language computes at least the functions that Turing machines
  compute, since it is Turing-complete, and at most these, since a Turing
  machine can run its interpreter. Theorem~\ref{thm:ceasefire} is the
  case of the $\lambda$-calculus and Turing's machines.
\end{proof}

\begin{remark}
  The ceasefire was signed in 1937, eleven years before the first
  stored-program computer ran a program~\cite{baby}. It concerns which
  functions can be computed, not how quickly, how clearly, or how often
  incorrectly. The rest of this paper is about those.
\end{remark}

\newpage % start the references on a new page
\begin{thebibliography}{99}
  \bibitem{church} A.~Church, An unsolvable problem of elementary number
  theory, \emph{American Journal of Mathematics} 58(2), 1936.
  \bibitem{kleene} S.~C. Kleene, $\lambda$-definability and recursiveness,
  \emph{Duke Mathematical Journal} 2(2), 1936.
  \bibitem{sondergaard} H.~S{\o}ndergaard and P.~Sestoft, Referential
  transparency, definiteness and unfoldability, \emph{Acta Informatica}
  27(6), 1990.
  \bibitem{strachey} C.~Strachey, Fundamental concepts in programming
  languages, \emph{Higher-Order and Symbolic Computation} 13(1--2), 2000.
  Lecture notes from 1967.
  \bibitem{turing36} A.~M. Turing, On computable numbers, with an
  application to the Entscheidungsproblem, \emph{Proceedings of the London
  Mathematical Society} s2-42(1), 1936.
  \bibitem{turing37} A.~M. Turing, Computability and $\lambda$-definability,
  \emph{Journal of Symbolic Logic} 2(4), 1937.
  \bibitem{baby} University of Manchester, How a 70-year-old Baby changed
  the face of modern computing,
  \url{https://www.manchester.ac.uk/about/news/how-a-70-year-old-baby-changed-the-face-of-modern-computing/},
  21 June 2018.
\end{thebibliography}

\end{document}
